\documentclass[12pt,a4paper]{article}
\usepackage[utf8]{inputenc}
\usepackage[russian]{babel}
\usepackage{geometry}
\usepackage{listings}
\usepackage{xcolor}
\usepackage{hyperref}
\usepackage{graphicx}
\usepackage{amsmath}
\usepackage{algorithm}
\usepackage{algpseudocode}

\geometry{margin=2.5cm}

% Настройка подсветки синтаксиса
\lstset{
    language=C++,
    basicstyle=\ttfamily\small,
    keywordstyle=\color{blue}\bfseries,
    commentstyle=\color{green!60!black},
    stringstyle=\color{red},
    numberstyle=\tiny\color{gray},
    breaklines=true,
    frame=single,
    numbers=left,
    stepnumber=1,
    tabsize=4,
    showspaces=false,
    showstringspaces=false
}

\title{Руководство программиста\\Компилятор языка my\_lang}
\author{Версия 1.0}
\date{\today}

\begin{document}

\maketitle
\tableofcontents
\newpage

\section{Введение}

Данное руководство описывает архитектуру компилятора языка \texttt{my\_lang} и предназначено для разработчиков, желающих расширить или модифицировать компилятор. Компилятор реализован на C++20 и использует автоматическую генерацию парсера и лексера на основе описания грамматики.

\section{Архитектура компилятора}

Компилятор состоит из следующих основных компонентов:

\begin{enumerate}
    \item \textbf{Лексер (Lexer)} --- разбивает исходный код на токены
    \item \textbf{Парсер (Parser)} --- строит абстрактное синтаксическое дерево (AST)
    \item \textbf{Семантический анализатор} --- проверяет корректность программы
    \item \textbf{Оптимизатор AST} --- выполняет простые оптимизации
    \item \textbf{Генератор POLIZ} --- создает промежуточное представление в виде обратной польской записи
\end{enumerate}

\subsection{Структура проекта}

\begin{verbatim}
project/
├── main.cpp                    # Точка входа
├── CMakeLists.txt              # Файл сборки CMake
├── grammar.txt                 # Описание грамматики
├── lexer_generator.py          # Генератор лексера
├── parser_generator.py         # Генератор парсера
├── system_reserved_identifiers.txt  # Зарезервированные слова
└── out/                        # Сгенерированные файлы
    ├── lexer.hpp              # Заголовок лексера
    ├── parser.hpp             # Заголовок парсера
    ├── parser.cpp             # Реализация парсера
    ├── ast.hpp                # Определения AST
    ├── ast.cpp                # Реализация AST
    ├── ast_utils.hpp          # Утилиты AST
    ├── ast_utils.cpp          # Реализация утилит
    ├── poliz.hpp              # Заголовок генератора POLIZ
    ├── poliz.cpp              # Реализация генератора POLIZ
    └── keywords.hpp           # Определения токенов
\end{verbatim}

\section{Лексер}

\subsection{Обзор}

Лексер реализован в файле \texttt{out/lexer.hpp} и использует алгоритм Ахо-Корасика для распознавания ключевых слов. Лексер разбивает исходный код на токены следующих типов:

\begin{itemize}
    \item Ключевые слова (\texttt{int}, \texttt{if}, \texttt{while}, и т.д.)
    \item Идентификаторы
    \item Числовые литералы
    \item Строковые литералы
    \item Символы и операторы
\end{itemize}

\subsection{Класс Lexer}

Основной класс лексера:

\begin{lstlisting}
namespace lexer {
    class Lexer {
    public:
        Lexer(const std::string& text, 
              const std::string& keywords_file);
        Token next();
        std::vector<Token> get_all_tokens();
    };
}
\end{lstlisting}

\subsection{Структура Token}

Каждый токен содержит следующую информацию:

\begin{lstlisting}
struct Token {
    parser::TokenType type;  // Тип токена
    std::string name;        // Имя токена
    std::string value;       // Значение токена
    size_t line;             // Номер строки
    size_t col;              // Номер столбца
};
\end{lstlisting}

\subsection{Алгоритм Ахо-Корасика}

Лексер использует алгоритм Бора для эффективного распознавания ключевых слов. Алгоритм строит конечный автомат, который позволяет находить все вхождения ключевых слов за один проход по тексту.

\subsection{Обработка комментариев}

Лексер автоматически пропускает:
\begin{itemize}
    \item Однострочные комментарии (\texttt{// ...})
    \item Многострочные комментарии (\texttt{/* ... */})
    \item Пробельные символы
\end{itemize}

\section{Парсер}

\subsection{Обзор}

Парсер реализован в файлах \texttt{out/parser.hpp} и \texttt{out/parser.cpp}. Парсер использует рекурсивный спуск для построения AST на основе грамматики, описанной в файле \texttt{grammar.txt}.

\subsection{Генерация парсера}

Парсер автоматически генерируется скриптом \texttt{parser\_generator.py} на основе файла \texttt{grammar.txt}. Грамматика описывается в формате BNF (Backus-Naur Form).

\subsection{Основные функции парсера}

Главная функция парсера:

\begin{lstlisting}
namespace parser {
    ast::AstNode* parse(lexer::Lexer& lexer, 
                        const std::string& path);
}
\end{lstlisting}

Для каждого нетерминала грамматики генерируется функция вида \texttt{TOK\_<NONTERMINAL>()}, которая возвращает соответствующий узел AST.

\subsection{Обработка ошибок}

Парсер использует исключения для обработки ошибок:

\begin{lstlisting}
struct ParseError : public std::runtime_error {
    using std::runtime_error::runtime_error;
};
\end{lstlisting}

При обнаружении синтаксической ошибки выбрасывается исключение \texttt{ParseError} с описанием проблемы.

\section{Абстрактное синтаксическое дерево (AST)}

\subsection{Обзор}

AST реализовано в файлах \texttt{out/ast.hpp} и \texttt{out/ast.cpp}. Дерево состоит из узлов различных типов, соответствующих конструкциям языка.

\subsection{Базовый класс AstNode}

Все узлы AST наследуются от базового класса:

\begin{lstlisting}
class AstNode {
public:
    NodeType type;
    int line;
    int col;
    std::vector<AstNode*> children;
    
    virtual ~AstNode();
    virtual void print(int depth = 0) const;
    virtual std::string to_string() const;
    void add_child(AstNode* child);
};
\end{lstlisting}

\subsection{Типы узлов}

Узлы AST делятся на два основных типа:

\begin{enumerate}
    \item \textbf{Терминальные узлы (TerminalNode)} --- представляют токены (идентификаторы, литералы, операторы)
    \item \textbf{Нетерминальные узлы} --- представляют синтаксические конструкции (выражения, операторы, объявления)
\end{enumerate}

\subsection{Примеры узлов}

\begin{itemize}
    \item \texttt{TOK\_PROGRAMNode} --- корневой узел программы
    \item \texttt{TOK\_DECLARATIONNode} --- объявление переменной или функции
    \item \texttt{TOK\_EXPRESSIONNode} --- выражение
    \item \texttt{TOK\_IFSTMTNode} --- условный оператор
    \item \texttt{TOK\_WHILESTMTNode} --- цикл while
\end{itemize}

\section{Семантический анализ}

\subsection{Обзор}

Семантический анализ реализован в файле \texttt{out/ast\_utils.cpp}. Функция \texttt{semantic\_check} выполняет проверку корректности программы на уровне семантики.

\subsection{Основная функция}

\begin{lstlisting}
namespace ast {
    bool semantic_check(AstNode* root);
}
\end{lstlisting}

Функция возвращает \texttt{true}, если семантические проверки пройдены успешно, и \texttt{false} в противном случае.

\subsection{Типы проверок}

Семантический анализатор выполняет следующие проверки:

\begin{itemize}
    \item Проверка типов в выражениях
    \item Проверка существования переменных и функций
    \item Проверка корректности использования операторов
    \item Проверка области видимости
\end{itemize}

\section{Оптимизация AST}

\subsection{Обзор}

Оптимизация AST реализована в функции \texttt{optimize\_ast} в файле \texttt{out/ast\_utils.cpp}.

\subsection{Типы оптимизаций}

В текущей реализации выполняются следующие оптимизации:

\begin{itemize}
    \item \textbf{Constant folding} --- вычисление константных выражений на этапе компиляции
    \item Простые локальные замены
\end{itemize}

\subsection{Пример оптимизации}

Выражение \texttt{2 + 3} может быть заменено на \texttt{5} на этапе компиляции.

\section{Генерация POLIZ}

\subsection{Обзор}

Генератор POLIZ реализован в файлах \texttt{out/poliz.hpp} и \texttt{out/poliz.cpp}. POLIZ (Польская Инверсная Запись) --- это промежуточное представление программы, удобное для интерпретации стековой машиной.

\subsection{Класс Poliz}

\begin{lstlisting}
namespace poliz {
    class Poliz {
    private:
        std::vector<PolizItem> items;
    public:
        void generate(ast::AstNode* node);
        const std::vector<PolizItem>& get_items() const;
        void print();
    };
}
\end{lstlisting}

\subsection{Типы операций}

Генератор POLIZ поддерживает следующие типы операций:

\begin{itemize}
    \item \texttt{PUSH\_VAL} --- поместить значение в стек
    \item \texttt{ADD}, \texttt{SUB}, \texttt{MUL}, \texttt{DIV}, \texttt{MOD} --- арифметические операции
    \item \texttt{ASSIGN} --- присваивание
    \item \texttt{EQ}, \texttt{NEQ}, \texttt{LT}, \texttt{GT}, \texttt{LEQ}, \texttt{GEQ} --- операции сравнения
    \item \texttt{GOTO} --- безусловный переход
    \item \texttt{JMP\_FALSE} --- условный переход
    \item \texttt{PRINT} --- вывод
    \item \texttt{INPUT} --- ввод
    \item \texttt{MEMBER} --- доступ к полю объекта
    \item \texttt{CALL} --- вызов функции
\end{itemize}

\subsection{Структура PolizItem}

\begin{lstlisting}
struct PolizItem {
    OpType op;              // Тип операции
    std::string value;      // Значение (для операндов)
    int jump_index;         // Индекс перехода (для GOTO)
};
\end{lstlisting}

\section{Генераторы кода}

\subsection{Генератор парсера}

Скрипт \texttt{parser\_generator.py} читает файл \texttt{grammar.txt} и генерирует:
\begin{itemize}
    \item Функции парсера для каждого нетерминала
    \item Определения узлов AST
    \item Заголовочные файлы
\end{itemize}

\subsection{Генератор лексера}

Скрипт \texttt{lexer\_generator.py} читает файл \texttt{system\_reserved\_identifiers.txt} и генерирует код лексера с поддержкой всех зарезервированных слов.

\section{Расширение компилятора}

\subsection{Добавление нового ключевого слова}

Для добавления нового ключевого слова:

\begin{enumerate}
    \item Добавьте ключевое слово в файл \texttt{system\_reserved\_identifiers.txt}
    \item Добавьте соответствующий тип токена в \texttt{keywords.hpp} (или перегенерируйте)
    \item Обновите грамматику в \texttt{grammar.txt}, если необходимо
    \item Пересоберите проект
\end{enumerate}

\subsection{Добавление нового оператора}

Для добавления нового оператора:

\begin{enumerate}
    \item Добавьте токен оператора в грамматику (\texttt{grammar.txt})
    \item Обновите правила парсинга выражений
    \item Добавьте обработку оператора в генератор POLIZ
    \item Пересоберите проект
\end{enumerate}

\subsection{Добавление нового типа данных}

Для добавления нового типа данных:

\begin{enumerate}
    \item Добавьте ключевое слово типа в \texttt{system\_reserved\_identifiers.txt}
    \item Обновите грамматику для поддержки нового типа
    \item Добавьте обработку типа в семантическом анализаторе
    \item Обновите генератор POLIZ для поддержки операций с новым типом
\end{enumerate}

\subsection{Добавление новой синтаксической конструкции}

Для добавления новой синтаксической конструкции:

\begin{enumerate}
    \item Добавьте правило в грамматику (\texttt{grammar.txt})
    \item Перегенерируйте парсер
    \item Добавьте соответствующий узел AST (если необходимо)
    \item Реализуйте обработку конструкции в семантическом анализаторе
    \item Реализуйте генерацию POLIZ для новой конструкции
\end{enumerate}

\section{Отладка}

\subsection{Вывод AST}

Для отладки можно использовать функцию печати AST:

\begin{lstlisting}
ast::print_tree(root, 0);
\end{lstlisting}

\subsection{Вывод POLIZ}

Для просмотра сгенерированного POLIZ:

\begin{lstlisting}
poliz_generator.print();
\end{lstlisting}

\subsection{Логирование}

В коде парсера и лексера можно добавить отладочный вывод для отслеживания процесса разбора.

\section{Тестирование}

\subsection{Структура тестов}

Тесты находятся в файле \texttt{autotest/autotest\_semantic.cpp} и используют фреймворк GoogleTest.

\subsection{Запуск тестов}

Для запуска тестов:

\begin{lstlisting}[language=bash]
cd cmake-build-debug
ctest
\end{lstlisting}

\subsection{Написание новых тестов}

Пример теста:

\begin{lstlisting}
TEST(ParserTest, SimpleExpression) {
    std::string source = "int x = 5 + 3;";
    lexer::Lexer lex(source, "../out/system_reserved_identifiers.txt");
    ast::AstNode* root = parser::parse(lex, "test.txt");
    ASSERT_NE(root, nullptr);
    // Проверки...
    delete root;
}
\end{lstlisting}

\section{Производительность}

\subsection{Оптимизации}

Текущая реализация использует следующие оптимизации:

\begin{itemize}
    \item Алгоритм Ахо-Корасика для быстрого распознавания ключевых слов
    \item Рекурсивный спуск для парсинга (эффективен для LL-грамматик)
    \item Простые оптимизации AST для уменьшения размера дерева
\end{itemize}

\subsection{Ограничения производительности}

\begin{itemize}
    \item Парсер не кэширует результаты разбора
    \item AST не оптимизируется агрессивно
    \item Генерация POLIZ выполняется за один проход без оптимизаций
\end{itemize}

\section{Будущие улучшения}

Возможные направления развития:

\begin{itemize}
    \item Реализация интерпретатора POLIZ
    \item Добавление более сложных оптимизаций AST
    \item Поддержка модулей и импорта
    \item Генерация машинного кода вместо POLIZ
    \item Улучшение сообщений об ошибках
    \item Поддержка отладки
\end{itemize}

\section{Заключение}

Данное руководство описывает архитектуру компилятора и способы его расширения. Для получения дополнительной информации об использовании компилятора обратитесь к руководству пользователя.

При возникновении вопросов или проблем рекомендуется изучить исходный код соответствующих компонентов и использовать отладочные инструменты для диагностики.

\end{document}

